% `tex.web` §307 and §444: two things an error context gets wrong together
% whenever the error is raised inside an output routine.
%
% §307 lists the token-list types and §314 prints their names. The routine is
% `output_text` and shows as `<output> `; pushed as an ordinary backed-up list
% it shows as `<to be read again> `, which is what ours did — so every error
% raised inside `\output`, of which trip.tex has many, named the wrong kind of
% level while quoting exactly the right text.
%
% §444 absorbs the space that ends a number rather than backing it up:
%
%     if cur_cmd<>spacer then back_input
%
% so `loc` sits past it and the read half of the context ends with the space.
% Trimming it moved the break one column left on every context that follows a
% number scan — `\penalty-10000` where TeX has `\penalty-10000 ` — and in
% trip.log that is most of them. The unread half *is* trimmed, because §31
% strips a line's trailing blanks as it reads it, so TeX has none to print.
%
% Both are visible only from inside an output routine that fails, because that
% is where the file position is examined between the number scan and whatever
% reads next.
\catcode`\{=1 \catcode`\}=2
\vsize=50pt
\output={\undefinedyy\shipout\box255}
\hbox{}\penalty-10000 % some trailing text on the line, so the break is visible
\hbox{}
\message{PROBE:done}
\end
